\subsubsection{Análisis de la Interpretación de Fechas y Razonamiento Cronológico}

En este experimento se evalúa la capacidad de los modelos de lenguaje GPT-3.5 y GPT-4 para interpretar fechas y realizar razonamiento cronológico. El objetivo es determinar si estos modelos pueden entender y calcular correctamente la edad de una persona en base a fechas y datos previamente proporcionados. El código completo de este análisis está disponible en el repositorio\footnote{\url{https://github.com/aditzend/hyperpersonalization/blob/main/app/notebooks/8.ipynb}}.

\paragraph{Configuración del Entorno}

Se utilizan FAISS para la búsqueda semántica y OpenAI Embeddings para la vectorización de mensajes. Se define una función para interactuar con los modelos GPT-3.5 y GPT-4 mediante la API de OpenAI.

La transcripción íntegra de este listado figura en el \textit{Anexo técnico} (véase el Apéndice~\ref{anexo:code-4-1-8-1}).

\paragraph{Creación de Contextos Conversacionales}

Se crean varios contextos que contienen conversaciones simuladas entre el usuario y el sistema. En uno de los contextos, el usuario solicita un turno para su madre, llamada Alba, y menciona que tiene 66 años. Otro contexto incluye una pregunta sobre la capital de Colombia.

La transcripción íntegra de este listado figura en el \textit{Anexo técnico} (véase el Apéndice~\ref{anexo:code-4-1-8-2}).

\paragraph{Vectorización y Almacenamiento}

Los mensajes de estos contextos se vectorizan utilizando embeddings y se almacenan en FAISS junto con sus respectivos metadatos, lo que permite realizar búsquedas y recuperar información relacionada.

La transcripción íntegra de este listado figura en el \textit{Anexo técnico} (véase el Apéndice~\ref{anexo:code-4-1-8-3}).

\paragraph{Consulta Simple sobre la Edad}

Se realiza una consulta simple: ``¿Qué edad tiene Alba?''. FAISS encuentra el contexto relevante, y el modelo GPT-3.5 responde correctamente que ``Alba tiene 66 años'', basándose en la información previamente proporcionada.

La transcripción íntegra de este listado figura en el \textit{Anexo técnico} (véase el Apéndice~\ref{anexo:code-4-1-8-4}).

El resultado obtenido fue: ``Alba tiene 66 años.''

\paragraph{Razonamiento Cronológico con GPT-3.5}

A continuación, se plantea una pregunta más compleja: ``¿Qué edad tenía Alba en 2020?''. Aquí se espera que el modelo entienda que si Alba tiene 66 años en 2023 (año actual de acuerdo a la fecha de realización del experimento), entonces en 2020 tendría 63 años. Sin embargo, GPT-3.5 falla en su respuesta.

La transcripción íntegra de este listado figura en el \textit{Anexo técnico} (véase el Apéndice~\ref{anexo:code-4-1-8-5}).

GPT-3.5 responde incorrectamente: ``En el contexto de la conversación anterior, Alba tenía 66 años en 2020.'', indicando que el modelo no ajusta la edad de acuerdo con la fecha solicitada.

\paragraph{Prueba con GPT-4}

Se realiza la misma consulta sobre la edad de Alba en 2020, esta vez utilizando GPT-4.

La transcripción íntegra de este listado figura en el \textit{Anexo técnico} (véase el Apéndice~\ref{anexo:code-4-1-8-6}).

GPT-4 responde: ``Alba nació en 1966. Por lo tanto, en 2020, tenía 54 años.'' Aunque GPT-4 realiza el cálculo correctamente, comete un error inicial al interpretar la edad como el año de nacimiento. A pesar de corregir el cálculo, este error inicial demuestra una falla en la comprensión de los datos proporcionados.

\paragraph{Conclusiones de la octava iteración}

Tanto GPT-3.5 como GPT-4 presentan dificultades significativas al interpretar fechas y realizar razonamientos cronológicos basados en la edad de una persona. Mientras que GPT-3.5 no realiza el cálculo de la edad correctamente, GPT-4 comete un error de interpretación antes de realizar el cálculo adecuado.

Los resultados demuestran que:

\begin{itemize}
    \item Los modelos pueden manejar consultas simples sobre datos proporcionados directamente
    \item Presentan serias dificultades cuando se les pide razonar cronológicamente
    \item GPT-3.5 no ajusta la edad según diferentes fechas
    \item GPT-4 malinterpreta inicialmente los datos numéricos (edad vs.\ año de nacimiento)
\end{itemize}

Esto sugiere que ambos modelos necesitan mejoras en la interpretación de datos numéricos y cronológicos cuando se extraen de contextos conversacionales. Es necesario implementar mecanismos adicionales para extraer y procesar datos numéricos correctamente en aplicaciones de hiperpersonalización. 